\sscp{Subgrouping implications from these patterns of distribution}{15-04-05}
\ili{As} can be seen from perusing the distribution maps in the previous
subsections and throughout this work,
there are many words or features that are found somewhere \ili{east}
of the Blust line,
but not west of it, or vice versa.
Some are found on both sides of the line.
Many of the features that begin to appear somewhere \ili{east} of the
line are not replacement innovations.
There are no features yet identified that precisely align with the Blust line.
Nor does the distribution of one feature align with the distribution of other features.
So it is the boundaries of several different Wallacean subgroups
which loosely share the presence or absence of certain features
that seem to define the Blust line (excepting \tsc{Taliabu}).
But patterns of distribution in only some subgroups or only some
nodes within subgroups indicate that many of those features have
been spread through contact (with Papuan languages or with other
Austronesian languages), and from diffusion, and thus have no subgrouping value of themselves.
Many do not seem to be patterned in systematic ways in eastern
Indonesia until one reaches the \tsc{Oceanic} subgroup.

\ili{As} one travels from the west to the \ili{east},
it is not completely clear when one enters the zone of linguistic Wallacea.
It depends entirely on which features are under scrutiny.
There are multiple thresholds,
some of which partially align with each other,
and many of which do not.

In the previous sections we have looked at a variety of overlapping patterns of distribution.
They include: 

\begin{enumerate}
	\item	Replacement innovations (e.g.
				PMP *mata ‘eye’ replaced with \tsc{Sula-Buru} **dama ‘eye’);
	\item	Words that are found in both Austronesian
				and Papuan languages, with a focus on those that show widespread
				consistency; \item	Words that are found on both sides of the Blust line,
				but should not be assigned to PMP;
	\item	Words that are found within Linguistic Wallacea,
				but only in a few subgroups;
	\item	Words that are found in several Wallacean subgroups,
				but only in certain nodes within those subgroups;
	\item	Words that appear to be exclusive to linguistic Wallacea
				extending out to Oceania.
				These do not appear to be found west of the Blust line.
				None have yet been found that are confidently reflected in all Wallacean subgroups.
\end{enumerate}

Pattern {\#}1 provides good supporting evidence for a subgroup
that is also based on more reliable evidence.
Patterns 2--5 all indicate that diffusion
and contact must be part of the consideration.
The sixth is the distributional pattern that is behind the CEMP hypothesis.
The fourth and fifth have contributed to the CMP hypothesis.
However, patterns 2--6 together beg the question,
“Why should the distribution of a word like **malip ‘laugh,
smile’ be given any more weight than the distribution of other
words that are greater than or less than CEMP,
but are not supported by exclusive innovations in the historical
phonology and morphosyntax?” The answer is simple,
“They shouldn’t be.” 
So in retrospect, it appears that it was the \it{assumption}
of the existence of a subgroup bounded at the Blust line that
makes the distribution of **malip ‘laugh,
smile’ and perhaps a few other words appear to support the CEMP hypothesis.
Methodologically, however, unless it aligns with stronger evidence
from the historical phonology
and morphosyntax, it is no better
and no worse than the other distributional patterns or isoglosses.
We need to be careful about letting a subgrouping hypothesis
and ideas of subgrouping boundaries determine the significance
of the distribution of lexical data.
All data are not equal,
and lexical data (which is easily borrowed) is the weakest evidence for subgrouping.

Words like {\#}muku ‘banana’ provide independent evidence that
words like **malip ‘laugh,
smile’, and **maya ‘tongue’, can disperse unevenly through our
target region without necessarily coming from a common Austronesian
parent language, nor having to be replacement innovations.
The distribution of shared developments in the lexicon within
the CEMP region does not require a proto language (PCEMP).

Because the overlapping distribution of these sorts of words
does not align with any subgroup defined by exclusively shared
innovations in the historical phonology
and morphosyntax, we must look for other explanations for their
uneven distribution in Linguistic Wallacea.